\documentclass[11pt]{article}
\usepackage{amsmath,amssymb}
\usepackage{xurl}
\usepackage[hidelinks]{hyperref}
\setlength{\emergencystretch}{2em}
% Shared commands for notes in docs/paper-gaps/.
%
% Two halves.  The link commands below are project identity: OWNER/REPO is the
% placeholder that scripts/bootstrap_project.py replaces with this project's
% GitHub slug and Pages site.  Everything after them is NOTATION, and notation
% belongs to the source paper: the definitions here are an example set, and a
% project replaces them with the macros of its own paper's style file.  The one
% rule that never changes: a command defined here must mean exactly what the
% source paper's macro means, or a note silently says something else.

\newcommand{\Lean}{\textsc{Lean}}
\newcommand{\leanid}[1]{\textsf{#1}}
\newcommand{\ghissue}[1]{\href{https://github.com/OWNER/REPO/issues/#1}{GitHub issue \##1}}
\newcommand{\ghpr}[1]{\href{https://github.com/OWNER/REPO/pull/#1}{PR \##1}}
% Local issue tracker (issues/NNNN-*.md in this repository); no remote URL.
\newcommand{\localissue}[1]{issue \##1 (\path{issues/})}
\newcommand{\doclink}[1]{\href{https://OWNER.github.io/REPO/docs/#1.html}{\path{#1}}}

% --- Notation: replace with the source paper's own macros --------------------
% \F, \N, \C, \R, \tr, \ot, \eps, \E, \bx, \bu, \bv, \by

\newcommand{\F}{\mathbb{F}}
\newcommand{\N}{\mathbb{N}}
\newcommand{\C}{\mathbb{C}}
\newcommand{\R}{\mathbb{R}}
\newcommand{\tr}{\mathrm{tr}}
\newcommand{\eps}{\epsilon}
\newcommand{\ot}{\otimes}
\newcommand{\E}{\mathop{\bf E\/}}

% bold random variables
\newcommand{\bu}{\boldsymbol{u}}
\newcommand{\bv}{\boldsymbol{v}}
\newcommand{\bx}{{\boldsymbol{x}}}
\newcommand{\by}{\boldsymbol{y}}

% T1 so literal < and > in \texttt placeholders typeset as themselves
\usepackage[T1]{fontenc}

% bra-ket
\usepackage{braket}

% --- Example structured spaces ---
\newcommand{\polyfunc}[3]{\mathcal{P}(#1, #2, #3)}
\newcommand{\polymeas}[3]{\mathrm{PolyMeas}(#1, #2, #3)}
\newcommand{\polysub}[3]{\mathrm{PolySub}(#1, #2, #3)}

% Approximation relations (paper uses \approx_\eps and \simeq_\zeta)
\newcommand{\approxeps}{\approx_{\varepsilon}}
\newcommand{\simgeq}{\simeq_{\zeta}}

\renewcommand{\ghissue}[1]{%
  \href{https://github.com/Dengnifer/MIPStarRE-QPBT/issues/#1}{GitHub issue \##1}}
\providecommand{\gapnote}[2]{\par\noindent\textbf{Verdict:} #1 (#2).\par}

\title{Complex Overlaps of the Constructed Subline Measurement}
\author{}
\date{2026-09-22}

\begin{document}
\maketitle
\gapnote{formalization-deviation}{open-proof}

\section{At a glance}

\textbf{Difficulty.} The constructed-measurement domain correction and the
complex-modulus estimates for all three scalar claims on the directly indexed
subline law are formalized. The source distribution correspondence remains open.
\textbf{Estimated weight.} The completed scalar proof combines the concrete
X marginal, expanded line-point consistency, restricted-distribution
averaging, and complex Cauchy--Schwarz. Positivity on opposite registers proves
reality of the Z overlap. The pasting consistency construction supplies the
line error for the concrete measurement. Remaining transport work concerns
the extended-line carrier and its evaluations, not this scalar calculation.
\textbf{Mathlib/project split.} Finite sums, positivity, and real
square roots use Mathlib; the measurement construction and subline law use
the local quantum API. \textbf{Key mathematical inputs.} Completeness of a
POVM, projectivity, and Cauchy--Schwarz suffice for the counterexample and
the scalar proof. The counterexample below is a mathematical
calculation, not a kernel-certified Lean instance.

\section{Key theorem forms}

Let $\widehat\psi$ be the normalized expanded state of a projective strategy
winning the Pauli basis test with probability at least $1-\eps$. Write
$L^W_f(\ell)$ and $M^W_a(u)$ for its expanded line and point effects. For
the subline law $D$, abbreviate the two correlations by
\begin{align*}
 A(T)&=\E_D\E_{(x,z,\alpha,\beta)\in\ell}
   \sum_{f_X,f_Z}\langle\widehat\psi,
   (T^{\ell_X,\ell_Z}_{f_X,f_Z}\ot
     M^Z_{f_Z(z)}(z)M^X_{f_X(x)}(x))\widehat\psi\rangle,\\
 B(T)&=\E_D\E_{(x,z,\alpha,\beta)\in\ell}
   \sum_{f_X,f_Z}\langle\widehat\psi,
   (T^{\ell_X,\ell_Z}_{f_X,f_Z}\ot
     M^Z_{f_Z(z)}(z))\widehat\psi\rangle.
\end{align*}
The paper's Claim 17-2 states
\[
 |A(T)-B(T)|\leq O\bigl(m\sqrt{\delta_{\mathrm{Line}}(\eps)}\bigr)
 \quad\text{for}\quad
 T^{\ell_X,\ell_Z}_{f_X,f_Z}
 =L^X_{f_X}(\ell_X)L^Z_{f_Z}(\ell_Z)L^X_{f_X}(\ell_X).
\]
Let $A_{\rm dir}$ and $B_{\rm dir}$ denote the same expressions for the
auxiliary subline law whose coordinate indices are sampled uniformly and
directly, with completed evaluations. The established Lean result is
\[
 |A_{\rm dir}(T)-B_{\rm dir}(T)|
 \le C m\sqrt{\delta_{\mathrm{Line}}(\eps)},
\]
for a universal $C>0$, the concrete sandwich above, and every directly
indexed subline witness. It assumes no joint-point or paired-line
consistency witness. This is not yet the estimate on the source law.
Let $J_{\rm dir}(T)$ denote the same overlap with the joint point effect
$Q^{x,z}_{f_X(x),f_Z(z)}$ in place of the ordered product. The other two
directly indexed complex estimates are
\begin{align*}
 |J_{\rm dir}(T)-A_{\rm dir}(T)|&\leq 2\sqrt{\delta_Q},\\
 B_{\rm dir}(T)&\in\mathbb R,\\
 |B_{\rm dir}(T)-1|&\leq
 C_Z\sqrt m\bigl(\delta_P^{1/4}+\delta_Q^{1/4}+\eps^{1/4}\bigr).
\end{align*}
Here the joint point family has the self-consistency and both ordered-product
bounds of the combined-points lemma. The first estimate needs no line
consistency assumption. For the third, the existing real-part proof has
$C_Z=2$ and uses the paired-line consistency error $\delta_P$. For a
polynomially controlled point-error function, the existing pasting theorem
produces a polynomial line-error function for the defined X-Z-X measurement;
the concrete specialization constructs its line witness internally.
No identity with an unspecified existential line witness is assumed.

The former Lean declaration quantified instead over every paired-line POVM
with the degree support and evaluated consistency conclusions of the
combined-lines lemma. Its scalar errors $\delta_P,\delta_Q$ were arbitrary,
independently of $\eps$. The right-hand side remained the same, with
$\delta_{\mathrm{Line}}(\eps)=\sqrt\eps$.

\section{The construction used by the source}

The paired-line measurement is defined by the displayed sandwich in the
proof of the combined-lines lemma. Claim 17-2 uses this definition to sum
over the Z polynomial:
\[
 \sum_{f_Z}T^{\ell_X,\ell_Z}_{f_X,f_Z}
 =L^X_{f_X}(\ell_X)
   \Bigl(\sum_{f_Z}L^Z_{f_Z}(\ell_Z)\Bigr)L^X_{f_X}(\ell_X)
 =L^X_{f_X}(\ell_X).
\]
The last equality uses completeness and projectivity. Cauchy--Schwarz
reduces the scalar difference to the square root of the resulting
X-line/point consistency deficit. The separate X marginal of the subline
law is sufficient; no joint conditional law for $(x,z)$ is needed.
Restriction to a line kind and coordinate costs a factor $2m$, so the
resulting square-root estimate is bounded by the paper's stated error
because $m\geq1$.\footnote{The construction is at lines 942--949 and the
claim and calculation at lines 1168--1201 of
\path{references/qpbt-paper/14_analysis_of_the_pauli_basis_test.tex},
labels \path{lem:qld-xz-lines} and \path{claim:17-2} in
\cite{Ji2020MIPStar}. The repair is tracked by \ghissue{414}; the earlier
proof attempt is \ghissue{405}.}

\section{A counterexample to the enlarged domain}

Choose an admissible tuple with $q>1$ and a perfect projective strategy.
The honest strategy in the completeness lemma supplies one; for example,
$(q,m,d)=(8,1,1)$ is admissible. Its expanded state has norm one. For
$a\in\F_q$, let $c_a$ be the constant polynomial with value $a$. For both
players and every line pair, define a POVM by
\[
 T_{c_a,c_b}=q^{-2}I\quad(a,b\in\F_q),\qquad
 T_{f_X,f_Z}=0\quad\text{for every other pair}.
\]
There are $q^2$ nonzero effects, so they sum to $I$. Both axis-degree
requirements hold. At any incident point a constant polynomial has a
uniquely determined value, even on a zero-direction line. Thus the optional
evaluation in Lean returns the ordinary field value, never an undefined
outcome, for every nonzero effect in this example.

The evaluated POVM is uniform on the $q^2$ point pairs. Against any complete
joint point measurement $Q$, its agreement is $q^{-2}$ and its consistency
defect is $1-q^{-2}$. Therefore the line-witness requirements hold with
$\delta_P=1$ in every opposite placement. A point witness also exists within
the former public domain: take the projective measurement which always
answers $(0,0)$. It is exactly self-consistent and has both ordered-product
distances at most $4$. Indeed, the squared norms of either ordered product
family sum to one, by projectivity and completeness, and
\begin{align*}
 \sum_{a,b}\|(Q_{a,b}-M^X_aM^Z_b)\widehat\psi\|^2
 &\leq 2\sum_{a,b}\|Q_{a,b}\widehat\psi\|^2\\
 &\quad+2\sum_{a,b}\|M^X_aM^Z_b\widehat\psi\|^2\\
 &=4.
\end{align*}
The reversed ordering satisfies the same calculation. Hence $\delta_Q=4$
is allowed by every point-witness field. These constants are permitted by
the former scalar quantifiers; the argument does not require them to be
polynomially small functions of $\eps$.

For each supported subline sample, completeness gives
\begin{align*}
 A(T)&=q^{-2}\sum_{a,b}\langle\widehat\psi,
    (I\ot M^Z_bM^X_a)\widehat\psi\rangle=q^{-2},\\
 B(T)&=q^{-2}\sum_{a,b}\langle\widehat\psi,
    (I\ot M^Z_b)\widehat\psi\rangle=q^{-1}.
\end{align*}
The first equality uses $\sum_bM^Z_b\sum_aM^X_a=I$, without a commutation
assumption. Probability averaging preserves these constants. The absolute
difference is therefore $(q-1)/q^2>0$, whereas the claimed bound is zero at
$\eps=0$. For $q=8$ the difference is $7/64$. The same perfect strategy is
admissible at arbitrarily small positive $\eps$, so excluding zero or
enlarging a universal constant cannot repair the claim. No polynomially
vanishing error depending on $\eps$ alone absorbs this difference.

\section{Blueprint and formal comparison}

The concrete measurement domain is retained. The scalar theorem now
compares complex expectations with the complex modulus, preserving the
universal constant, factor $m$, and square-root error. It is named
\leanid{subline\_remove\_X\_factor\_direct}. Its distribution is still the
directly indexed law represented by \leanid{SubLineWitness}, not the
source law $D$. Thus it is a formalization-only estimate, linked from
\path{lem:claim-17-2-direct}, and is not a formalization of the source
claim. The source statement remains at \path{lem:claim-17-2}, with neither
a Lean link nor a completion mark.

The same scalar distinction affects Claim 17-1: an ordered product
$M^Z_bM^X_a$ need not be Hermitian, so a bound on the real part of its
correlation does not bound its complex modulus. The proved real-part
estimate is retained as
\leanid{subline\_replace\_by\_ordered\_product\_re\_direct}.
Its bound is $2\sqrt{\delta_Q}$: the squared-distance comparison on
one placement is $4\delta_Q$ and uses $W^\dagger W$, not $W^2$.
The complex version now has the same bound, for the defined sandwich, as
\leanid{subline\_replace\_by\_ordered\_product\_direct}.
Likewise \leanid{subline\_Z\_term\_near\_one\_re\_direct} retains the
proved real Z-overlap estimate, and
\leanid{subline\_Z\_overlap\_eq\_real\_direct} identifies its complex
overlap with that real number. Thus
\leanid{subline\_Z\_term\_near\_one\_direct} has the same error bound
in complex modulus. The concrete specialization is
\leanid{subline\_concrete\_Z\_term\_near\_one\_direct}.
All these conclusions retain the directly indexed domain.
The three existing source blueprint statements are reused; no duplicate
source theorem with a new proof hole is introduced.

The distribution discrepancy is documented in
\path{docs/paper-gaps/qpbt_ld-dimension-divisibility.tex}. Admissibility
does not supply the divisibility needed to instantiate the source line law
at dimension $2m+2$. The current correction neither adds that divisibility
as a hypothesis nor identifies the directly indexed law with the source law.
Scalar realignment was tracked in \ghissue{474} (closed 2026-09-09) and
review of PR~478; the live tracker is \ghissue{118}.
The source-distribution work under \ghissue{118} remains open;
the direct complex estimates are recorded under \ghissue{689}.
Its historical budget is unchanged.

The real-part analogue of the second estimate is retained under its
explicitly qualified name \leanid{subline\_remove\_X\_factor\_re\_direct},
with auxiliary blueprint node \path{lem:claim-17-2-direct-real}. The
deprecated compatibility names \leanid{subline\_replace\_by\_ordered\_product},
\leanid{subline\_remove\_X\_factor} and \leanid{subline\_Z\_term\_near\_one}
keep the aliasing they have on the base branch.

The abstract combined-line witness record remains unchanged. It records a
complete paired-line POVM, the two axis-degree conditions, and consistency on
every directed opposite placement. Its consistency field, however, compares
the completed evaluation alphabet obtained by adjoining an undefined outcome,
whereas Lemma~\path{lem:qld-xz-lines} sums only over field-valued evaluations.
Consequently the record captures the measurement and degree data used by the
source construction but does not yet express its field-valued consistency
conclusion. The source blueprint node therefore remains open. Both witness
existence proofs instantiate the record with the concrete sandwich. A later
argument needing that construction can use the consistency theorem and the
same explicit record; merely choosing an unspecified existential witness does
not recover its defining formula.\footnote{The formalization-only theorem is
\leanid{combined\_line\_measurement\_consistency} in
\path{MIPStarRE/QPBT/Combining/Lines/Construction.lean}. Its hypotheses explicitly supply
a polynomially controlled point family. Its proof was recovered for
\ghissue{512} from the saved point-to-line comparisons and pasting argument,
including restoration of discarded zero-direction probability. The source-facing
\leanid{exists\_combinedLinesWitness} has an unchanged statement.}

The concrete scalar calculation is now proved: the exact X marginal,
expanded line-point consistency, and the separate restricted X-point mixture
give the deficit bound, and Cauchy--Schwarz closes both the real-part and
complex-modulus estimates.
No marginal identity, equality, or producer was added to the public
hypotheses. In particular, this proof does not use the
paired-line consistency theorem.\footnote{The deficit theorem is
\leanid{exists\_concreteXPointOverlap\_deficit\_le} in
\path{MIPStarRE/QPBT/Combining/Lines/ConcreteXDeficit.lean}; the scalar
closure is tracked by \ghissue{480}.}

\section{Source correspondence of the three scalar claims}

All three current scalar theorems use the directly indexed auxiliary
subline law, not the seed-indexed law $D$ of the source. The latter requires
an extended-dimension index function not defined for every admissible
parameter tuple, as detailed in the dimension-divisibility note.\footnote{See
\path{docs/paper-gaps/qpbt_ld-dimension-divisibility.tex}.}
The proposed construction target
\leanid{subline\_source\_distribution\_transport} must establish the
carrier, evaluation, and separate point-marginal identities wherever the
source law is defined. Its unrestricted use still requires resolution of
that definition-level obstruction; no divisibility assumption is added to
the paper theorem.

All three scalar comparisons now have complex-modulus versions on the auxiliary
law. The ordered product $M_ZM_X$ need not be Hermitian. With $W=Q-M_ZM_X$,
the squared norm in the first claim therefore uses $W^\dagger W$, not the
printed $W^2$. This corrects the proof calculation without changing the claimed
asymptotic error or assuming commutation of $M_X,M_Z$.
After regrouping the line POVM by its two completed evaluations, write its
effects as $A_a$ and its right-hand differences as $W_a$. Complex
Cauchy--Schwarz for the positive sesquilinear forms
$\langle u,A_av\rangle$ gives
\begin{align*}
 \left|\mathbb E\sum_a\langle\widehat\psi,A_aW_a\widehat\psi\rangle\right|
 &\leq
 \left(\mathbb E\sum_a\langle\widehat\psi,A_a\widehat\psi\rangle\right)^{1/2}\\
 &\quad\cdot\left(\mathbb E\sum_a
   \langle W_a\widehat\psi,A_aW_a\widehat\psi\rangle\right)^{1/2}\\
 &\leq \left(\mathbb E\sum_a\|W_a\widehat\psi\|^2\right)^{1/2}.
\end{align*}
The first factor is one by completeness and normalization; the second inequality
uses $0\leq A_a\leq I$. The completed point families have zero effects at
undefined outcomes. Their distance sum consequently reduces exactly to the
field-valued sum. The joint uniform-point marginal and the proved same-placement
distance bound give $4\delta_Q$, hence $2\sqrt{\delta_Q}$.
This discharges the former complex-overlap target without a reality premise.
The proof reuses the positive-operator square-root construction and complex
inner-product Cauchy--Schwarz from Mathlib through the existing weighted
overlap lemma.\footnote{The new distance estimate and its subline specialization
are in \path{MIPStarRE/QPBT/Combining/SubLineComplex.lean}.}
For the second claim, the existing complex weighted estimate uses positivity,
completeness, and projectivity.\footnote{This is
\leanid{norm\_overlap\_gap\_le\_sqrt\_one\_sub\_of\_isProj} in
\path{MIPStarRE/QPBT/Combining/ComplexOverlapGap.lean}.}
For the third claim, the two opposite-placement positive factors commute.
Their product is positive, so its expectation has zero imaginary part.
Finite sums and averages with real weights preserve this equality. Mathlib's
characterization of a positive semidefinite matrix by its quadratic form
proves this fact; it is not an assumed reality condition. It follows that
$|B_{\rm dir}-1|$ is exactly $|\operatorname{Re}B_{\rm dir}-1|$,
with no error loss.

The blueprint retains the three proved real-part auxiliary estimates and
the complex second estimate. The five new auxiliary entries supporting the
complex estimates on the directly indexed law carry statement and proof
completion marks. The three source-labelled claims remain uncertified.
For the same concrete measurement the triangle inequality now gives
\[
 |J_{\rm dir}(T)-1|\leq 2\sqrt{\delta_Q}
 +Cm\sqrt{\delta_{\mathrm{Line}}(\eps)}
 +C_Z\sqrt m\bigl(\delta_P^{1/4}+\delta_Q^{1/4}+\eps^{1/4}\bigr).
\]
This supplies the complex scalar calculation in the first route of the
combined-line proof on the auxiliary law. It does not supply the source law,
field-valued evaluation transport, the stronger joint conditional law used by
the second route, or the printed error form for the extended-line theorem.
The complete local consumer comparison is in the dedicated audit
\path{audits/2026-09-22-subline-complex-estimates.md}.

The recovered theorem \leanid{combined\_line\_measurement\_consistency}
proves consistency of the same unconditioned X-Z-X measurement
with the supplied polynomially controlled joint point family on the
original product of line-point distributions, in every opposite placement.
The construction and its discharge by the source point-to-line comparison and
pasting argument are described in
\path{docs/paper-gaps/qpbt_combined-lines-error-term.tex}.

\section{Statement integrity after merging the scalar estimates}

The paper assumes the projective setting, the constructed X-Z-X measurement,
and the source subline law; its conclusion bounds the complex overlap
difference by $C m\sqrt{\delta_{\mathrm{Line}}(\eps)}$.
The second scalar theorem retains its existing projective-setting and
directly indexed subline hypotheses,
the defined measurement, and the same complex-modulus conclusion and error
scale. Its proof introduces no public assumptions. The real-part analogue
of the second estimate and the first and third real-part estimates retain
their statements. The new first estimate keeps the actual X-Z-X measurement
and the joint point bounds, drops the unused line-consistency witness, and
proves the complex modulus with constant $2$. The new Z-reality identity holds
for any line POVM, including the concrete sandwich. The complex third bound
inherits exactly the real theorem's assumptions and error, while its concrete
specialization derives line consistency internally from the polynomial point
error already present in the source construction.
The verdict relative to the source remains a scope restriction to the
directly indexed law with optional evaluations; the complex scalar
conclusions for all three estimates are now proved on that domain.

\section{Conclusion}

The discrepancy lies in the formal quantification over arbitrary witnesses,
not in the paper's constructed-measurement assertion. The counterexample
forces a domain correction. Using the defined sandwich restores the measurement
domain. Complex Cauchy--Schwarz and positivity close the three scalar
comparisons for the directly indexed auxiliary law;
transport to the source distribution remains
open. The source claims themselves retain their complex-modulus statements
without asserting proofs of their source-law forms.

\bibliographystyle{alpha}
\bibliography{references}
\end{document}
